% =============================================================================
% System architecture trade-off: localization technique (v2, first-principles sizing)
% Needs in the preamble: booktabs, tabularx, longtable, amsmath, xcolor[table]
%   \definecolor{passc}{RGB}{214,239,214}  \definecolor{failc}{RGB}{248,215,213}
% Tables in generated/ are written by tradeoff_model.py; re-run it after any
% input or requirement change and re-upload generated/*.tex to Overleaf.
% =============================================================================

\section{Concept Exploration and System Trade-off --- \textit{[Name]}}\label{sec:tradeoff}

\subsection{Trade Tree and Method}

Section~1.4 keeps the localization architecture open between a compact coded-mask
imager and a timing constellation, and names three priorities for the trade: cost per
usefully localized burst, end-of-life detector stability and compliance with the ESA
5-year disposal rule. Because the localization technique fixes the number of spacecraft,
the payload mass and the alert chain, it is traded first. The trades that follow from it
are listed in Table~\ref{tab:tradetree} and are carried out in the subsystem chapters.

\begin{table}[htbp]\centering\small
\caption{Trade tree. Each lower-level trade is opened only after the trade above it is closed.}\label{tab:tradetree}
\begin{tabularx}{\linewidth}{@{}l X X l@{}}\toprule
ID & Trade & Alternatives & Section \\\midrule
T-1 & Localization architecture & coded mask; timing triangulation; count-rate ratios; lobster-eye & \ref{sec:tradeoff} \\
T-2 & Orbit & inclination $0^\circ$--$30^\circ$ (ENV-050); altitude vs.\ 5-year lifetime (LIF-010) & Mission analysis \\
T-3 & Alert link & ground VHF network; commercial relay constellation; own relay & Telecom., ground segment \\
T-4 & Detector & CdTe; CdZnTe; scintillator + SiPM & Payload \\
T-5 & Launch & rideshare; dedicated small launcher & Launch \\\bottomrule
\end{tabularx}\end{table}

The trade follows the rules of the course report guide: every criterion is derived from a
requirement and carries a unit, every value is computed or cited, and the only qualitative
criterion (in-orbit evidence of the localization technique) is identified as such. Each
alternative is sized from physical relations against the SRR thresholds instead of being
scaled from an existing instrument, so that the result follows from the requirements and
not from the design choices of another mission.

\textbf{Stage 1 (screening)} rejects techniques that cannot meet a threshold requirement on
physical grounds. \textbf{Stage 2} sizes each remaining technique and evaluates it. The
decision rule is: (i) reject any alternative that misses a threshold of MIS-020, MIS-030,
MIS-040, PER-065 or the nominal figure of merit $N_u\ge30\,\mathrm{yr^{-1}}$ of
Section~2.3; (ii) among the rest, select the lowest cost per usefully localized burst.

No weighted scoring table is used. The report guide accepts weights only when they are
taken from the literature or from expert input, and no published weighting exists for GRB
localization architectures. The figure of merit of Section~2.3 combines the three
priorities of Section~1.4 in one number traceable to SN-001 and SN-009.

\subsection{Stage 1: Screening}

Table~\ref{tab:screening} shows the four techniques that have been flown or designed for
GRB localization. Count-rate ratios are rejected because the systematic error of
Fermi/GBM, the best-calibrated instrument of this type, is $3.71^\circ$
\cite{Connaughton2015}, more than twenty times the 10\,arcmin of MIS-030. Lobster-eye
optics are rejected because they work in the 0.5--4\,keV band \cite{Yuan2022} and do not
cover the 5--50\,keV band of PER-040.

% generated by tradeoff_model.py - do not edit by hand
\begin{table}[htbp]\centering\small
\caption{Stage-1 screening of localization techniques against the threshold requirements MIS-030 and PER-040. Values from \cite{Barthelmy2005,Connaughton2015,Fiore2020,Yuan2022}.}\label{tab:screening}
\begin{tabularx}{\linewidth}{@{}l l l l c X@{}}\toprule
Technique & Flown / designed on & Localization & Band [keV] & Go & Reason \\\midrule
Coded-aperture mask & Swift/BAT, INTEGRAL, SVOM & 1--4$'$ (BAT) & 15--150 (BAT) & \cellcolor{passc}Go & Arcminute localization flown; carried to stage 2 \\
Timing triangulation & Vela; HERMES-TP/SP, CAMELOT (design) & 2.5$^\circ$ (HP, pred.); 15$'$/coord. (goal) & 3--2000 & \cellcolor{passc}Go & Could meet MIS-030 with enough spacecraft; carried to stage 2 \\
Count-rate ratios & Fermi/GBM, BATSE, BurstCube & 3.71$^\circ$ systematic core & 8--40\,000 & \cellcolor{failc}No-go & Systematic floor $\gg$ 10$'$: fails MIS-030 \\
Lobster-eye focusing & Einstein Probe/WXT & arcmin (soft X-ray) & 0.5--4 & \cellcolor{failc}No-go & Band outside PER-040 (5--50 keV): fails PER-040 \\
\bottomrule\end{tabularx}\end{table}


\subsection{Stage 2: Alternatives}

\begin{itemize}
\item \textbf{A1, coded-mask imager sized to the PER-040 threshold.} One spacecraft; the
smallest camera that meets PER-040, PER-010/020 and PER-065.
\item \textbf{A2, coded-mask imager sized to the PER-040 goal}
($1\times10^{-8}$\,erg\,cm$^{-2}$\,s$^{-1}$). Shows whether a larger camera lowers the cost
per burst.
\item \textbf{A0, coded-mask imager limited to PHY-010.} The largest camera with a 10\,kg
payload.
\item \textbf{B1, timing-triangulation constellation.} HERMES-type 3U CubeSats
\cite{Ghirlanda2024,Fiore2020} in two orbital planes (ARC-070), with as many spacecraft as
needed for $r_{90}=10'$ at the most favourable published timing accuracy
\cite{Werner2018}.
\item \textbf{B0, HERMES Pathfinder as designed}, six units \cite{Ghirlanda2024}.
\end{itemize}

\subsection{Coded-Mask Sizing Model}

A coded-mask camera consists of a position-sensitive detector of side $D$ and area $A=D^2$,
a mask of the same size at a distance $f$ above it, and a passive shield closing the space
between them \cite{Braga2019}. The camera is sized in four steps, each tied to an SRR
threshold.

\paragraph{Flux requirement.} PER-040 gives the 1-s peak energy flux $F$ in the band
$E_1$--$E_2=5$--50\,keV. For the photon spectrum $N(E)\propto E^{-1}$, which corresponds to
the typical low-energy index of GRB spectra \cite{Band1993}, the photon flux is
\begin{equation}\label{eq:pflux}
P = \frac{F}{E_2-E_1}\,\ln\frac{E_2}{E_1},
\end{equation}
so the PER-040 threshold $F=2.5\times10^{-8}$\,erg\,cm$^{-2}$\,s$^{-1}$ corresponds to
$P_{\rm req}=0.80$\,ph\,cm$^{-2}$\,s$^{-1}$.

\paragraph{Detector area.} For a mask with open fraction $\rho$ decoded by balanced
correlation, the imaging signal-to-noise ratio of a source of photon flux $P$ is
\begin{equation}\label{eq:snr}
\mathrm{SNR} = \frac{\sqrt{\rho(1-\rho)}\;\varepsilon P\sqrt{A\,t}}{\sqrt{b+(1-\rho)\,\varepsilon P}},
\end{equation}
where $\varepsilon$ is the detector efficiency, $t$ the integration time and $b$ the
background count rate per unit detector area. Eq.~\eqref{eq:snr} follows from propagating
Poisson noise through the correlation decoding; $\rho=0.5$ maximizes it in the
background-dominated case \cite{Braga2019}. With $t=1$\,s (PER-040), $\mathrm{SNR}=6.5$
(PER-100), $b=3.0$\,counts\,cm$^{-2}$\,s$^{-1}$ (ENV-010) and $\varepsilon=0.8$, solving
Eq.~\eqref{eq:snr} for $A$ gives the detector area.

\paragraph{Field of view and resolution.} PER-065 requires
$\Omega\,\eta_{\rm eff}\ge0.75$\,sr. With $\eta_{\rm eff}=0.74$ from the team's orbit
simulation of a 535\,km, $28.5^\circ$ orbit (mission-analysis chapter), the half-coded field
of view must be $\Omega\ge1.01$\,sr. For a square field with half-angle $\alpha$ per axis,
\begin{equation}\label{eq:fov}
\Omega = 4\arcsin\!\left(\sin^2\alpha\right),\qquad f=\frac{D}{2\tan\alpha},
\end{equation}
Eq.~\eqref{eq:fov} gives $\alpha=30.0^\circ$ and fixes $f$ for a given $D$. The angular resolution
(FWHM) of the camera and the 90\,\% error radius are \cite{Braga2019}
\begin{equation}\label{eq:res}
\theta_{\rm res}=\frac{\sqrt{m^2+d^2}}{f},\qquad
r_{90}=\sqrt{\left(\frac{k_{\rm loc}\,\theta_{\rm res}}{\mathrm{SNR}}\right)^2+\sigma_{\rm sys}^2},
\end{equation}
where $m$ is the mask element size, $d$ the detector pixel pitch, $k_{\rm loc}=1.5$ the
ratio between the 90\,\% error radius and $\theta_{\rm res}/\mathrm{SNR}$, and
$\sigma_{\rm sys}=2.5'$ the systematic allocation of PER-020. Requiring $r_{90}=10'$ at
SNR~8 (PER-010) gives $\theta_{\rm res}\le51.6'$. With $m=1.25\,d$, the ratio used on
Swift/BAT \cite{Barthelmy2005}, Eq.~\eqref{eq:res} gives the pixel pitch $d$.

\paragraph{Mass and power.} With material densities and thicknesses from
Table~\ref{tab:tradeinputs}, the payload mass is
\begin{equation}\label{eq:mass}
m_{\rm pl} = (1+\mu)\Big[k_{\rm fee}\,\varrho_{\rm CdTe}\,t_{\rm det}A
 + k_{\rm supp}(1-\rho)\,\varrho_{\rm W}\,t_{\rm mask}A
 + \sigma_{\rm sh}\left(4Df+D^2\right) + m_{\rm DPU}\Big],
\end{equation}
where the four terms are the detector modules (crystal mass times the factor $k_{\rm fee}$
for front-end electronics and mechanics), the tungsten mask with its support frame, the
graded side and rear shield of areal density $\sigma_{\rm sh}$, and the data processing
unit. $\mu=0.2$ is the margin for a new development \cite{ESAmargin}. The payload power is
$N_{\rm pix}\,p_{\rm ch}+p_{\rm DPU}$ with $N_{\rm pix}=(D/d)^2$.

\paragraph{Plausibility of the background.} ENV-010 limits the background to
3.0\,counts\,cm$^{-2}$\,s$^{-1}$. The cosmic X-ray background alone, integrated over
5--50\,keV with the fit of \cite{Gruber1999} and seen through the open half of the mask
over $\Omega$, gives 2.3\,counts\,cm$^{-2}$\,s$^{-1}$. ENV-010 therefore leaves only
0.7\,counts\,cm$^{-2}$\,s$^{-1}$ for albedo, shield leakage and activation, and the case of
twice the ENV-010 background is included in the sensitivity analysis.

\paragraph{Detection and usefully localized bursts.} The SRR derives MIS-020 from PER-040 and
PER-065: a camera exactly at both thresholds detects $N_{\rm req}=40\,\mathrm{yr^{-1}}$.
A camera with flux limit $P_{\rm lim}$ and exposure $\Omega\,\eta_{\rm eff}$ detects
\begin{equation}\label{eq:ndet}
N_{\rm det}=N_{\rm req}\left(\frac{P_{\rm lim}}{P_{\rm req}}\right)^{-\beta}
\frac{\Omega\,\eta_{\rm eff}}{0.75\,\mathrm{sr}},
\end{equation}
where $\beta$ is the slope of the integral burst distribution $N(>P)\propto P^{-\beta}$;
$\beta=1$ is assumed and varied between 0.8 and the Euclidean limit 1.5. The flow-down in
Eq.~\eqref{eq:ndet} is
taken from the SRR and must be verified with the burst-population Monte Carlo of the
mission-analysis chapter. Of the triggers above $\mathrm{SNR}_{\rm a}=6.5$, those above
$\mathrm{SNR}_{10}=8$ reach $r_{90}\le10'$, so the usefully localized rate is
\begin{equation}\label{eq:nu}
N_u = N_{\rm det}\left(\frac{\mathrm{SNR}_{10}}{\mathrm{SNR}_{\rm a}}\right)^{-\beta} f_{\rm lat},
\end{equation}
with $f_{\rm lat}=1$ if the alert reaches GCN within 20\,min; Eq.~\eqref{eq:nu} is the figure
of merit of Section~2.3 for a coded mask. The team's simulation of a
VHF alert network gives 92.5\,\% of alerts within 30\,s for the 535\,km orbit.

\subsection{Triangulation Model}

For timing triangulation the error per coordinate is \cite{Fiore2020}
\begin{equation}\label{eq:tri}
\sigma_\theta \approx 2.4^\circ\,\frac{\sigma_t}{\sqrt{N_{\rm vis}-3}},
\end{equation}
with $\sigma_t$ the combined delay and clock error in ms, $N_{\rm vis}$ the number of
spacecraft that see the burst and an average baseline of 7000\,km;
$r_{90}=\sqrt{2\ln10}\,\sigma_\theta=2.146\,\sigma_\theta$. $\sigma_t=0.34$\,ms is derived
from CAMELOT, about $20'$ with nine CubeSats \cite{Werner2018}, and $\sigma_t=1.8$\,ms from
HERMES Pathfinder, about $2.5^\circ$ with six units \cite{Fiore2020}. A spacecraft sees a
given burst with probability $f_{\rm see}=0.70\,\Omega_{\rm unit}/4\pi=0.29$, where 0.70 is
the sky fraction not blocked by the Earth at 550\,km and $\Omega_{\rm unit}=5.2$\,sr
\cite{Ghirlanda2024}. The fleet size follows from Eq.~\eqref{eq:tri} as
$N=N_{\rm vis}/f_{\rm see}$. The detection rate is that of HERMES Pathfinder,
$214\,\mathrm{yr^{-1}}$ \cite{Ghirlanda2024}, and half of the detected bursts are assumed
to reach $r_{90}\le10'$ because the fleet is sized for the median burst.

\subsection{Cost Model}

The single-spacecraft dry mass is $3.3\,m_{\rm pl}$, the average of
\cite[Table~10-6, p.~312]{Wertz1999}, plus the 20\,\% system margin of DES-010. Its cost
follows Eq.~\eqref{eq:sscm}, the Small Satellite Cost Model \cite[Table~20-6, p.~797]{Wertz1999}:
\begin{equation}\label{eq:sscm}
C_{\rm bus} = 781 + 26.1\,m_{\rm bus}^{1.261}\ \ [\mathrm{FY00\,k\$}],\qquad
C_{\rm tot} = 1.495\,(1+0.4)\,C_{\rm bus},
\end{equation}
where $m_{\rm bus}$ is the bus dry mass in kg (model range 20--400\,kg), 0.4 is the SSCM
payload-to-bus cost ratio and 1.495 adds integration and test, program level, ground support
equipment and launch and orbit operations support \cite[p.~797]{Wertz1999}. Costs are
escalated to 2024 dollars with the consumer price index \cite{BLSCPI}. The SSCM is not valid
for CubeSats, so the constellation is costed by analogy: the HERMES-SP project cost of
3.43\,MEUR for development and three flight units \cite{CORDIS821896} sets the first-unit
cost, and the learning curve of Eq.~\eqref{eq:learning}
\begin{equation}\label{eq:learning}
C_N = C_{\rm TFU}\,N^{B},\qquad B = 1-\frac{\ln(1/S)}{\ln 2},
\end{equation}
with $S=85\,\%$ for more than 50 units \cite[Eq.~20-5, p.~809]{Wertz1999} gives the
production cost, and 5\,M\$ is added for constellation development \cite{Garcia2021}.
Launch is priced at 350\,k\$ for the first 50\,kg and 7\,k\$/kg above
\cite{SpaceXRideshare}, one launch per orbital plane. Operations are excluded for both
architectures and are estimated in the cost chapter.

\subsection{Results}

Table~\ref{tab:maskdesign} gives the three coded-mask designs and
Table~\ref{tab:tradematrix} the trade-off matrix.

% generated by tradeoff_model.py - do not edit by hand
\begin{table}[htbp]\centering\small
\caption{First-principles coded-mask designs, Eqs.~\eqref{eq:pflux}--\eqref{eq:mass}. A1 and A2 are sized to the PER-040 threshold and goal; A0 is the largest camera within PHY-010.}\label{tab:maskdesign}
\begin{tabular}{@{}l l rrr@{}}\toprule
Parameter & Unit & A1 & A2 & A0 \\\midrule
Half-coded field of view $\Omega$ & sr & 1.01 & 1.01 & 1.01 \\
FoV half-angle per axis $\alpha$ & deg & 30.0 & 30.0 & 30.0 \\
Angular resolution $\theta_{\rm res}$ (FWHM) & arcmin & 51.6 & 51.6 & 51.6 \\
Detector side $D$ (= mask side) & cm & 37.1 & 90.0 & 22.0 \\
Detector area $A$ & cm$^2$ & 1\,375 & 8\,098 & 482 \\
Effective area $\rho\,\varepsilon\,A$ & cm$^2$ & 550 & 3\,239 & 193 \\
Mask--detector distance $f$ & cm & 32.1 & 77.8 & 19.0 \\
Pixel pitch $d$ & mm & 3.01 & 7.30 & 1.78 \\
Pixels & -- & 15\,197 & 15\,197 & 15\,197 \\
1-s flux limit at SNR 6.5 & $10^{-8}$\,erg\,cm$^{-2}$\,s$^{-1}$ & 2.50 & 1.00 & 4.37 \\
Detector modules & kg & 4.0 & 23.7 & 1.4 \\
Mask and support & kg & 2.0 & 11.7 & 0.7 \\
Shield & kg & 9.2 & 54.1 & 3.2 \\
Data processing unit & kg & 3.0 & 3.0 & 3.0 \\
Payload mass incl. 20\,\% margin & kg & 21.8 & 111.1 & 10.0 \\
Payload power & W & 12.6 & 12.6 & 12.6 \\
\bottomrule\end{tabular}\end{table}

% generated by tradeoff_model.py - do not edit by hand
\begin{table*}[htbp]\centering\footnotesize
\caption{Stage-2 trade-off matrix for the localization architecture. Green: threshold met; red: not met. PHY-010 and PHY-020 are open [TBC] values; they are reported but not used to reject. The last row is the only qualitative criterion; it is used because the in-orbit maturity of a localization chain cannot be expressed as a number. All values are computed by \texttt{tradeoff\_model.py} (Appendix~\ref{app:tradeoff}).}\label{tab:tradematrix}
\begin{tabularx}{\linewidth}{@{}l l l l XXXXX@{}}\toprule
Criterion & Unit & Thr. & Req. & A1: mask at PER-040 threshold & A2: mask at PER-040 goal & A0: mask within PHY-010 & B1: triangulation, sized & B0: HERMES-class, 6 s/c \\\midrule
Localization $r_{90}$ at SNR 8 & arcmin & $\le$10 & MIS-030 & \cellcolor{passc}10.0 & \cellcolor{passc}10.0 & \cellcolor{passc}10.0 & \cellcolor{passc}10.0 & \cellcolor{failc}556.2 \\
GRBs detected & yr$^{-1}$ & $\ge$40 & MIS-020 & \cellcolor{passc}40 & \cellcolor{passc}100 & \cellcolor{failc}23 & \cellcolor{passc}214 & \cellcolor{passc}214 \\
$\Omega\,\eta_{\rm eff}$ & sr & $\ge$0.75 & PER-065 & \cellcolor{passc}0.75 & \cellcolor{passc}0.75 & \cellcolor{passc}0.75 & \cellcolor{passc}fleet & \cellcolor{passc}fleet \\
Alerts within 30\,s & \% & $\ge$66 & MIS-040 & \cellcolor{passc}92.5 & \cellcolor{passc}92.5 & \cellcolor{passc}92.5 & \cellcolor{failc}0.0 & \cellcolor{failc}0.0 \\
Usefully localized, $N_u$ & yr$^{-1}$ & $\ge$30 & FoM & \cellcolor{passc}32.5 & \cellcolor{passc}81.2 & \cellcolor{failc}18.6 & \cellcolor{failc}0.0 & \cellcolor{failc}0.0 \\
Payload mass (per s/c) & kg & $\le$10\,[TBC] & PHY-010 & \cellcolor{failc}21.8 & \cellcolor{failc}111.1 & \cellcolor{passc}10.0 & \cellcolor{passc}1.6 & \cellcolor{passc}1.6 \\
Payload power (per s/c) & W & $\le$15\,[TBC] & PHY-020 & \cellcolor{passc}12.6 & \cellcolor{passc}12.6 & \cellcolor{passc}12.6 & \cellcolor{passc}-- & \cellcolor{passc}-- \\
Spacecraft & -- & -- & -- & 1 & 1 & 1 & 394 & 6 \\
Launch mass (total) & kg & -- & -- & 87 & 440 & 40 & 2\,364 & 36 \\
Cost, launch incl., ops excl. & M\$$_{2024}$ & -- & -- & 23 & 155 & 10 & 185 & 12 \\
Cost per useful burst (2\,yr) & M\$/GRB & min & Decision & 0.35 & 0.95 & 0.28 & $\infty$ & $\infty$ \\
Localization evidence & qual. & -- & -- & Technique flown (Swift/BAT, INTEGRAL) & Technique flown (Swift/BAT, INTEGRAL) & Technique flown (Swift/BAT, INTEGRAL) & Not demonstrated in orbit & Not demonstrated in orbit \\
\midrule
All thresholds met & & & & \textbf{yes} & \textbf{yes} & no & no & no \\
\bottomrule\end{tabularx}\end{table*}


\textbf{A1} meets every threshold under the SRR flow-down of Eq.~\eqref{eq:ndet} (but see
Section~\ref{sec:flowdown}) with a 37\,cm square CdTe detector (1375\,cm$^2$,
550\,cm$^2$ effective area), a mask 32\,cm above it, 3.0\,mm pixels and 15\,200 channels.
The payload has a mass of 21.8\,kg including margin and draws 12.6\,W; the shield accounts
for more than half of the mass. The 87\,kg spacecraft costs about 23\,M\$, or 0.35\,M\$ per
usefully localized burst over two years. A1 is sized exactly at PER-040 and therefore
detects exactly 40 bursts per year: it carries no performance margin.

\textbf{A2} detects 100 bursts per year, but its 90\,cm detector and 111\,kg payload make the
spacecraft 440\,kg and 155\,M\$, and the cost per burst rises to 0.95\,M\$. Mass grows
faster than yield because the shield scales with $D\,f\propto A$ while the yield scales with
$A^{\beta/2}$.

\textbf{A0} shows that PHY-010 and MIS-020 are incompatible for a coded mask: a 10\,kg
payload allows 482\,cm$^2$ and 23 detected bursts per year. The 10\,kg value of PHY-010 was
derived from the mass per area of a non-imaging HERMES detector (Appendix~B), which a mask
and shield cannot reach.

\textbf{B1} reaches $r_{90}=10'$ only with 394 CubeSats (2364\,kg, 185\,M\$), even with the
most favourable published timing accuracy. Its decisive weakness is latency: triangulation
needs the photon arrival times of at least four spacecraft, and HERMES plans to deliver
trigger messages within 30\,min and photon data within 24\,h \cite{Fiore2024}. This fails
MIS-040 and the 20\,min usefulness limit, so $N_u=0$. \textbf{B0} gives
$r_{90}\approx9^\circ$ and fails MIS-030 by a factor of more than 50.

\subsection{Sensitivity Analysis}

% generated by tradeoff_model.py - do not edit by hand
\begin{table*}[htbp]\centering\footnotesize
\caption{Sensitivity of the selection to every assumption and uncertain input (one at a time). B1$^*$ is B1 with on-board inter-satellite triangulation assumed, which no flown or funded mission has demonstrated. Cost per useful burst in M\$/GRB over 2 years.}\label{tab:sensitivity}
\begin{tabularx}{\linewidth}{@{}X r r r r r r c c@{}}\toprule
Case & A1 kg & A1 M\$/GRB & A2 kg & A2 M\$/GRB & B1 s/c & B1$^*$ M\$/GRB & Selected & Selected if ISL \\\midrule
Baseline & 21.8 & 0.35 & 111.1 & 0.95 & 394 & 0.87 & A1 & A1 \\
$k_{\rm loc}$ = 1.0 & 21.8 & 0.35 & 111.1 & 0.95 & 394 & 0.87 & A1 & A1 \\
$k_{\rm loc}$ = 2.0 & 21.8 & 0.35 & 111.1 & 0.95 & 394 & 0.87 & A1 & A1 \\
Background 2 $\times$ ENV-010 & 38.3 & 0.66 & 214.1 & 2.15 & 394 & 0.87 & A1 & A1 \\
Detector efficiency 0.6 & 35.3 & 0.60 & 192.7 & 1.88 & 394 & 0.87 & A1 & A1 \\
$\eta_{\rm eff}$ = 0.50 & 19.8 & 0.31 & 99.0 & 0.83 & 394 & 0.87 & A1 & A1 \\
Shield areal density 2 $\times$ & 32.9 & 0.55 & 176.0 & 1.68 & 394 & 0.87 & A1 & A1 \\
Detector module factor 4 & 24.7 & 0.40 & 128.1 & 1.13 & 394 & 0.87 & A1 & A1 \\
$\beta$ = 0.8 & 21.8 & 0.33 & 111.1 & 1.10 & 394 & 0.87 & A1 & A1 \\
$\beta$ = 1.5 (Euclidean) & 21.8 & 0.39 & 111.1 & 0.67 & 394 & 0.87 & A2 & A2 \\
Dry/payload mass ratio = 2 & 21.8 & 0.17 & 111.1 & 0.39 & 394 & 0.87 & A1 & A1 \\
Dry/payload mass ratio = 7 & 21.8 & 1.00 & 111.1 & 2.96 & 394 & 0.87 & A1 & B1* \\
Payload cost = 0.8 $\times$ bus & 21.8 & 0.45 & 111.1 & 1.22 & 394 & 0.87 & A1 & A1 \\
Triangulation $\sigma_t$ = 0.2 ms & 21.8 & 0.35 & 111.1 & 0.95 & 145 & 0.41 & A1 & A1 \\
CubeSat TFU halved & 21.8 & 0.35 & 111.1 & 0.95 & 394 & 0.48 & A1 & A1 \\
Launch price doubled & 21.8 & 0.35 & 111.1 & 0.97 & 394 & 0.94 & A1 & A1 \\
\bottomrule\end{tabularx}\end{table*}


Table~\ref{tab:sensitivity} varies every assumption one at a time. The coded mask remains the
selected architecture in every case; the payload mass of A1 ranges from 19.8 to 38.3\,kg, the
largest increases coming from a doubled background (38.3\,kg) and a lower detector efficiency
(35.3\,kg). The localization constant $k_{\rm loc}$ does not change the mass, because the
field of view fixes $f/D$, but it sets the pixel count: $k_{\rm loc}=2$ requires 2.3\,mm
pixels, 27\,000 channels and 18.5\,W, which exceeds PHY-020.

Two results matter for the baseline. First, for the Euclidean slope $\beta=1.5$, A1 falls to
$N_u=29.3\,\mathrm{yr^{-1}}$ and misses the figure of merit, and the larger A2 is selected:
a camera sized exactly at PER-040 has no margin. Second, the constellation would be preferred
only if on-board inter-satellite triangulation within 30\,s were possible (B1$^*$) and the
spacecraft dry mass were at the top of the SMAD range (seven times the payload mass). No
mission has demonstrated inter-satellite triangulation, so this is the condition under which
the trade should be reopened.

\subsection{Verification of the Detection-Rate Flow-down}\label{sec:flowdown}

Equation~\eqref{eq:ndet} assumes, as the SRR does, that a camera exactly at PER-040 and
PER-065 detects the 40 bursts per year of MIS-020. This assumption was checked with a
Monte Carlo simulation.

\begin{itemize}
\item \textbf{Purpose and requirement verified:} MIS-020, through the flow-down
PER-040\,+\,PER-065\,$\rightarrow$\,MIS-020.
\item \textbf{Tool:} MATLAB R2025b, \texttt{Trade-off/verification/check\_detection\_rate.m}.
\item \textbf{Model:} the burst population of the team's mission-analysis suite, used
unchanged: 900 bursts per year over the whole sky above a 1-s peak flux of
0.25\,ph\,cm$^{-2}$\,s$^{-1}$ (15--150\,keV), a broken power-law flux distribution and Band
spectra. $4\times10^5$ bursts are drawn inside the field of view of each camera and detected
with Eq.~\eqref{eq:snr} in the 5--50\,keV band, including the off-axis cosine.
\item \textbf{Main simplification:} partial coding inside the field of view is ignored, which
makes the detection rates optimistic.
\item \textbf{Plausibility check:} the same population, seen by a 1.4\,sr camera with a
threshold of 0.4\,ph\,cm$^{-2}$\,s$^{-1}$, gives 67 bursts per year before duty-cycle losses,
consistent with the Swift/BAT on-board rate of 63 per year quoted in Section~2.3.
\end{itemize}

\begin{table}[htbp]\centering\small
\caption{Detection rates of the three coded-mask designs in the burst-population Monte Carlo,
compared with the SRR flow-down of Eq.~\eqref{eq:ndet}.}\label{tab:flowdown}
\begin{tabular}{@{}l r r r r r@{}}\toprule
Design & $A$ [cm$^2$] & $N_{\rm det}$, Eq.~\eqref{eq:ndet} [yr$^{-1}$] & $N_{\rm det}$, MC [yr$^{-1}$] & $N_u$, MC [yr$^{-1}$] & $\beta_{\rm eff}$ \\\midrule
A1 & 1375 & 40.0  & \cellcolor{failc}23.5 & \cellcolor{failc}18.7 & 1.06 \\
A2 & 8098 & 100.0 & \cellcolor{passc}48.6$^\dagger$ & \cellcolor{passc}44.2$^\dagger$ & 0.27 \\
A0 & 482  & 22.9  & \cellcolor{failc}13.0 & \cellcolor{failc}10.0 & 1.22 \\\bottomrule
\multicolumn{6}{@{}l}{\footnotesize $^\dagger$ Lower bound: A2 still detects a burst at the population floor (SNR 7.9).}
\end{tabular}\end{table}

Table~\ref{tab:flowdown} shows that the flow-down does not hold in the team's population. A1
detects 23.5 bursts per year, not 40. The fraction of usefully localized triggers (0.80) and
the local slope of the distribution ($\beta_{\rm eff}=1.06$) agree with the values assumed in
Eqs.~\eqref{eq:ndet}--\eqref{eq:nu}. The difference therefore comes from the normalization of
the flow-down itself. With $\Omega\,\eta_{\rm eff}=0.75$\,sr, only 53.7 bursts per year above
the population floor fall in the field of view at all, so MIS-020 requires detecting three
quarters of them. Meeting MIS-020 needs either a larger field of view than PER-065 asks for or a
camera close to A2. PER-040 and PER-065 are consequently not sufficient for MIS-020 and must be
revised together with the camera sizing in the payload chapter. The architecture decision is
unaffected: the constellation fails MIS-040 independently of the detection rate.

\subsection{Selection and Consequences for the Baseline}

\textbf{A single spacecraft carrying a coded-mask imager is selected.} It is the only
architecture that meets all threshold requirements with a localization technique already
flown, and it has the lowest cost per usefully localized burst. The constellation is rejected
because it fails MIS-040 with any demonstrated alert chain and needs 394 spacecraft to meet
MIS-030.

The camera size is not yet fixed. Under the SRR flow-down, A1 is the cheapest compliant design;
the burst-population check of Section~\ref{sec:flowdown} shows that A1 detects only 23.5
bursts per year, so the payload chapter must size the area and the field of view together
against MIS-020 with the population model. The selection changes the requirement baseline as
follows:
\begin{itemize}
\item \textbf{PHY-010} (payload mass, [TBC]) is raised from 10\,kg: A1 alone needs 21.8\,kg
including margin, and the final value follows from the camera sizing.
\item \textbf{PER-040 and PER-065} are revised together, since meeting both thresholds does not
give the 40 bursts per year of MIS-020.
\item \textbf{ENV-010} (background) is confirmed by a background model in the payload chapter,
since the cosmic X-ray background alone uses 76\,\% of it.
\item \textbf{PER-010} is rewritten without the instrument-specific error law: ``$r_{90}\le10'$
for bursts with SNR\,$\ge$\,8''.
\item \textbf{Table~C}: ARC-050 to ARC-070 (constellation) are deleted.
\item The 87\,kg spacecraft sets the launch-mass class for T-5, and the $i\le30^\circ$ orbit
of ENV-050 must be checked against the inclinations rideshare services offer (T-2, T-5).
\end{itemize}

\paragraph{Limitations.} The mass model resolves only the main payload elements and relies on
assumed overhead factors, which the sensitivity analysis brackets. The detection rate relies on
the SRR flow-down from PER-040 to MIS-020 and on an assumed burst-distribution slope. The SSCM
dates from 1999 and its bus-cost relation has a standard error of 3.7\,M\$ (FY00)
\cite[p.~797]{Wertz1999}. Operations costs are excluded, and $\sigma_t$ for triangulation is
derived from predicted, not measured, performance.
